% Generated draft: requires mathematical and source review.
\documentclass[11pt]{article}
\usepackage[margin=25mm]{geometry}
\usepackage{fontspec}
\setmainfont{Latin Modern Roman}
\usepackage{amsmath,amssymb,mathtools,unicode-math}
\setmathfont{Latin Modern Math}
\usepackage{enumitem,needspace,xurl,hyperref}
\setlength{\parindent}{0pt}
\setlength{\parskip}{5pt}
\setlength{\emergencystretch}{3em}
\newcommand{\N}{\mathbb{N}}
\newcommand{\Z}{\mathbb{Z}}
\newcommand{\Q}{\mathbb{Q}}
\newcommand{\R}{\mathbb{R}}
\newcommand{\C}{\mathbb{C}}
\newcommand{\F}{\mathbb{F}}
\newcommand{\A}{\mathbb{A}}
\newcommand{\PP}{\mathbb{P}}
\newcommand{\E}{\mathbb{E}}
\newcommand{\eps}{\varepsilon}
\newcommand{\dd}{\,\mathrm{d}}
\begin{document}
\section*{UnsolvedMath ID: 700460}
\textbf{Source problem code:} catalogue entry
\subsection{Definitions and mathematical statement}
A group is cyclic if it is generated by a single element, finitely generated if it has a finite generating set, and finitely presented if it admits a presentation with finitely many generators and finitely many defining relations. For a group \(G\) and an element \(a\in G\), the condition that every element is conjugate to a power of \(a\) means
\[
\forall g\in G\;\exists h\in G\;\exists n\in\Z
\quad g=h a^n h^{-1}.
\]
Here powers include zero and negative integer powers, and both \(h\) and \(n\) may depend on \(g\).

\textbf{Part (b), the remaining question in the supplied record (D. V. Anosov).} Does there exist a noncyclic finitely presented group \(G\) and an element \(a\in G\) such that
\[
G=\bigcup_{n\in\Z}\{h a^n h^{-1}:h\in G\}?
\]

\textbf{Archived part (a) (D. V. Anosov).} Does there exist a noncyclic finitely generated group \(G\) containing an element \(a\) such that every element of \(G\) is conjugate to an integer power of \(a\)? The supplied archive reports an affirmative answer to this part, attributed to V. S. Guba. The registry labels the combined record \texttt{partially\_solved} and retains part (b) in the main section. These are supplied source and curation assessments, not independently verified mathematical or literature conclusions.
\subsection{Short English statement}
The remaining question asks whether a noncyclic group with a finite presentation can have every element conjugate to a power of one fixed element. The supplied archive reports an affirmative answer when finite presentation is replaced by finite generation.
\subsection{Sources}
\begin{itemize}
\item Supplied local UnsolvedMath registry record: permanent ID \(700460\), exact source code \texttt{catalogue entry}, title \emph{Kourovka Notebook Problem 8.8}. The record lists R. I. Grigorchuk as proposer and 1982 as the proposed year; both displayed parts are specifically attributed to D. V. Anosov.
\item \emph{The Kourovka Notebook (Unsolved Problems in Group Theory)}, 8th issue (1982), problem 8.8; supplied background reports an update in the 21st issue (2026), \texttt{arXiv:1401.0300v48}. The supplied collection description names E. I. Khukhro and V. D. Mazurov as editors.
\item Supplied Archive of Solved Problems attribution for the affirmative answer to part (a): V. S. Guba, \emph{Math. USSR--Izv.}, volume 29 (1986), pages 233--277. No article title is supplied.
\item Supplied status review dated \texttt{2026-10-16}: canonical status \texttt{partially\_solved}, evidence field \texttt{published}; the review assigns solved parts to the archive and remaining parts to the main section. This is supplied dated curation, not independently verified.
\end{itemize}
\end{document}
